\documentclass[../../main.tex]{subfiles}
\begin{document}

\begin{center}
  {\fontsize{48pt}{58pt}\selectfont \textbf{1995}\normalsize\textbf{年度}}
\end{center}

\begin{table}[htb]
  \centering
  \caption{出題テーマと難易度}
  \begin{tabular}{|c|p{9cm}|c|}
    \hline
    \textbf{問題番号} & \textbf{テーマ} & \textbf{難易度} \\
    \hline
    第1問 & 多項式の余りと二次方程式の根 & \\
    \hline
    第2問 & 円周上の弧の中点を繰り返す数列 & \\
    \hline
    第3問 & 二つの二次式の値の評価 & \\
    \hline
    第4問 & 行列の相似変換 & \\
    \hline
    第5問 & 札の交換ゲームと勝率 & \\
    \hline
    第6問 & 双曲線の接線と面積で定まる領域 & \\
    \hline
  \end{tabular}
\end{table}

\begin{table}[htb]
  \centering
  \caption{解答}
  \renewcommand{\arraystretch}{1.4}
  \begin{tabular}{|c|c|p{9cm}|}
    \hline
    \textbf{問題番号} & \textbf{小問} & \textbf{答え} \\
    \hline
    第1問 & (1) & 証明問題 \\
    \cline{2-3}
    & (2) & 範囲は解答の図を参照 \\
    \hline
    第2問 & (1)，(2) & 証明問題 \\
    \cline{2-3}
    & (3) & $\alpha_{3n}=\left(-\dfrac{1}{8}\right)^n\left(\alpha_0-\dfrac{2\pi}{3}\right)+\dfrac{2\pi}{3}$ \\
    \hline
    第3問 & -- & 証明問題 \\
    \hline
    第4問 & (1)，(2) &  \\
    \hline
    第5問 & (1) & $p_2=\dfrac{1}{4}$，$q_2=\dfrac{3}{4}$ \\
    \cline{2-3}
    & (2) & $(n+2)p_n-np_{n-2}=1$（証明問題） \\
    \cline{2-3}
    & (3) & $p_n=\dfrac{n}{2(n+2)}$（$n$偶数），$\dfrac{n+1}{2(n+2)}$（$n$奇数） \\
    \hline
    第6問 & (1) & $0<a<2$ \\
    \cline{2-3}
    & (2) & 範囲は解答の図を参照 \\
    \hline
  \end{tabular}
\end{table}

\paragraph{試験形態}
6問．試験時間は要確認．

\end{document}
